\section{Introduction}\label{sec:intro}

Alice keeps her coins in a two-key arrangement: her wallet holds one key and a custodian holds another, and both must sign. She wants three things from it. She wants no party, the custodian included, ever to hold the whole key. She wants to make one move when she spends and then go offline, while a service watches the fee market and finishes the job. And she wants the service's operators to co-sign without learning her address, her amount or her destination. Threshold signing gives the first, line dealing~\cite{LD} gives the second, and neither gives the third.

Line dealing starts from one observation. A Schnorr signature is the value at the challenge $c$ of the line $s(c) = k + c\,x$, and an ECDSA signature is the value at the sighash $h$ of the line $s(h) = h\,w + r\,\mu$ with $w = k^{-1}$ and $\mu = k^{-1}x$. The owner knows both coefficients and deals the line to a committee, either as finished signatures on a grid of conflicting transactions (sealed points) or as both coefficients under commitments that force their constant terms (open lines). In the open form the committee holds the key at threshold. In the sealed form the owner holds it alone.

\paragraph{The idea.} We cut the line. At address birth the key is split additively,
\[
  x = \xf + \tw + \sum_{a \in A} x_a ,
\]
into a \emph{committee half} $\xf$, which the committee holds in Shamir form, a secret \emph{tweak} $\tw$, and one \emph{shard} $x_a$ for each actor $a$ of a set $A$: the user, an issuer, a custodian, a guardian group, or any party the charter names. A Schnorr signature is linear in the key and in the nonce, so it splits the same way. Each actor computes its \emph{contribution}, a point of its own shard line at the challenge, and the committee adds its half with one linear partial per node. The committee half needs only a challenge scalar, a binding factor and a committee nonce. Every term that depends on the address or on an actor sits in the actors' contributions, which the committee never needs to read before it signs.

Three properties follow from this split, and they are the subject of the paper.
\begin{itemize}[leftmargin=1.5em]
\item \emph{Co-custody without a signing round.} Each actor signs offline from public nonce points, and the committee signs with one linear partial per node. ECDSA, which is not linear in the nonce, pays once: an interactive session at birth prepares one-use slots in which the committee and every actor hold additive halves of $k^{-1}$ and $k^{-1}x$.
\item \emph{The sealed guarantee for every actor.} A committee that is corrupt in full, together with a coalition of actors that the charter does not admit, holds its own half and every dealt contribution. It can complete what an admitted set sealed. Every other message needs a fresh point of a shard line that some honest actor never dealt. This claim is sketched on ECDSA and on Schnorr with an honest initiator, and its reduction for honest pool actors is open (\cref{sec:claims}).
\item \emph{Blind co-signing.} A node signs a \emph{ticket}: a few scalars and one aggregate actor nonce point. The ticket carries no address, no message and no actor key.
\end{itemize}

\paragraph{Contributions.} We give split-line dealing for BIP340, Ed25519 and ECDSA over nested access structures of AND, threshold and OR gates (\cref{sec:setting,sec:construction}). We give one envelope that carries every request type in one size. Its outer dealing deals one secret, and the committee derives the header, body and contribution keys from it with a threshold PRF whose partials leave a node only after the board orders the envelope. We give the slot rules that let the candidate tickets of a fee grid share one committee nonce per generation and input, the precedence rule that decides which request owns a slot, and the leaf-set binding that keeps two subsets of a group off one ECDSA slot. We state twenty-one security claims with their status (\cref{sec:security}): the algebra and the counting are proved in Lean~4, and each cryptographic claim is sketched or open with its named obligation. The Lean development has 57 theorems and one equivalence for split-line dealing, all on the standard axioms (\cref{sec:formal}). We state what each party sees (\cref{sec:blindness}) and measure the costs on a prototype (\cref{sec:eval}).

\paragraph{What is not new.} The two-nonce binding of each contribution is the binding of FROST~\cite{FROST,RFC9591}, with every actor and the committee as signers. Additive and Shamir sharing of a key, nested by gates, is the standard construction of monotone access structures. The additive ECDSA halves of a slot are presignatures in the sense of Groth and Shoup~\cite{GS21}, and the functionalities that prepare them have candidate instantiations in published threshold ECDSA protocols, whose adaptation at our threshold is open (\cref{sec:ecdsa-slots}). A user share that a network completes on a condition is the shape of the 2PC-MPC line~\cite{TwoPCMPC,FMMSS25,IkaDocs}. What we add is the cut of the dealt line into a committee half and outside shards, the sealed contributions that keep line dealing's confinement for every actor, the ticket that lets the committee sign without the address, and the slot and envelope rules that make one move enough.

\paragraph{Organization.} \Cref{sec:setting} fixes the parties, keys, access structures and profiles. \Cref{sec:construction} gives the protocol. \Cref{sec:security} states the claims and their status. \Cref{sec:blindness} states what each party sees. \Cref{sec:formal,sec:eval} report the Lean development and the measured costs. \Cref{sec:related} covers related work and \cref{sec:limits} the limits and open problems.
